\documentclass[a4paper,11pt,fleqn]{article}%fleqn pour aligner les equations à gauche
\usepackage{Cours}


\newcommand{\chnum}{Chapitre 2}
\newcommand{\chtitle}{Dosage par étalonnage}
\newcommand{\dtype}{TP}


\begin{document}
\maketitle

Le permanganate de potassium, ou cristal de Condy, est un composé ionique de formule \ce{KMnO4 (aq)}. Dissout dans l'eau, il forme une solution aqueuse violette d'ions \ce{K+(aq)} et \ce{MnO4-(aq)} qui est utilisée, à des concentrations différents, comme antiseptique cutanée, désinfectant pour laver les légumes, ou pour le nettoyage des aquariums en aquariophilie. Pour traiter l'eczéma, il est recommandé d'utiliser une préparation \SI{0,25}{g} dans \SI{2}{L} d'eau.\\
\textbf{Problématique : On dispose d'une solution $S$ de permanganate de potassium. Peut-on utiliser cette solution pour traiter l'eczéma ?}

\begin{figure}[h!]
    \begin{minipage}{.5\linewidth}
    \caption{\textbf{Liste du matériel}}
\begin{itemize*}
	\item solution $S'$ de permanganate de potassium
	\item 5 solutions étalons de permanganate de potassium de concentrations en masse $\gamma_i$ comprise entre \SI{10}{\milli\gram\per\liter} et \SI{50}{\milli\gram\per\liter}.
	\item solution étalon pour conductimètre
	\item conductimètre
	\item verrerie et appareils de mesure usuels
	\item ordinateur équipé de \emph{LatisPro} 
\end{itemize*}
\end{minipage}
 \begin{minipage}{.5\linewidth}
 \caption{\textbf{Informations indiquées sur le flacon de permanganate de potassium}}
$M$ = \SI{158,0}{\gram\per\mole}\\
\includegraphics[scale=.43]{Ressources/TSPE_C2_AE1.png}
\end{minipage}
 \begin{minipage}{.5\linewidth}
 \caption{\textbf{Ressources disponibles}}
\begin{itemize*}
	\item documents p. 50 du livre
	\item fiche méthode du conductimètre
\end{itemize*}
\end{minipage}
\end{figure}
\noindent

\section*{Dosage conductimétrique}
	\begin{enumerate}
	\item Déterminer s'il serait possible de réaliser un dosage conductimétrique pour déterminer la concentration en masse de permanganate de potassium dans la solution $S$. \textcolor{gray!80!blue}{\textbf{ANA, APP}}\\
	\end{enumerate}
La solution $S$ à étudier est trop concentrée pour être utilisable en conductimétrie. On la dilue donc d'un facteur 10 avant de l'analyser. On obtient alors une solution $S'$ de concentration en masse (ou titre massique) $\gamma '$. Celle-ci est déjà préparée.

	\begin{enumerate}[resume]
	\item Proposer une démarche expérimentale pour déterminer la concentration en masse de permanganate de potassium dans la solution $S$ à l'aide d'un dosage conductimétrique. \textcolor{gray!80!blue}{\textbf{ANA, COM}}\\
	\begin{center}
	\textbf{Appel}
	\end{center}
	\item Mettre en œuvre la démarche proposée. \textcolor{gray!80!blue}{\textbf{REA}}\\
	\end{enumerate}

	Pour l'utilisation du conductimètre, se référer à la fiche proposée et choisir le mode "sans compensation de température". Exploiter les données à l'aide de \textit{Latis Pro}. \textcolor{gray!80!blue}{\textbf{REA}}

	\begin{enumerate}[resume]
	\item Répondre à la problématique. \textcolor{gray!80!blue}{\textbf{ANA}}
	\item Vérifier que la gamme de concentrations de la gamme étalon est adaptée à l'étude. \textcolor{gray!80!blue}{\textbf{VAL}}
	\end{enumerate}

\end{document}